% Class-compatibility preamble retained for a reader excerpt of
%   Sahar Diskin, Anna Geisler, "Perfect Matching in Product Graphs and in their
%   Random Subgraphs", arXiv:2404.14020v4, CC BY 4.0.
%
% The source pins its own publisher class (scrartcl); the delivery wrapper accepts
% only the classes article and amsart, so this file supplies the amsart baseline and
% then copies the author's own declaration block mechanically and unchanged
% (pinned-original-source.tex, source lines 30-57: the author's newtheorem
% declarations and the author's macro declarations).  The author's own package
% list of source lines 4-21 is reused only for the entries the retained text needs
% (enumitem, amsmath, amssymb, amsthm, mathtools, xcolor, geometry) and the
% cross-reference packages already needed by the author's own \Cref/\eqref usage
% (hyperref, cleveref); tikz, caption, subcaption, float, graphicx, wrapfig and
% abstract serve figures and the title block that this excerpt does not carry.
% No mathematical macro, statement, symbol or assumption is introduced, changed
% or removed by this file.
\documentclass[10pt,a4paper]{amsart}
\usepackage[shortlabels]{enumitem}
\usepackage{amsmath}
\usepackage{amssymb}
\usepackage{amsthm}
\usepackage{mathtools}
\usepackage{xcolor}
\usepackage[margin=1in]{geometry}
\usepackage{hyperref}
\usepackage{cleveref}
\hypersetup{
  colorlinks=true,
  linkcolor=blue,
  filecolor=magenta,
  urlcolor=cyan
}

\newtheorem{theorem}{Theorem}
\newtheorem{thm}{Theorem}[section]
\newtheorem{corollary}[thm]{Corollary}
\newtheorem{lemma}[thm]{Lemma}
\newtheorem{proposition}[thm]{Proposition}
\newtheorem{claim}[thm]{Claim}
\newtheorem{remark}[thm]{Remark}
\newtheorem{question}[thm]{Question}

\newcommand{\whp}{whp}
\newcommand{\defi}{\operatorname{def}}
\newcommand{\odd}{\operatorname{odd}}
\newcommand{\iso}{\operatorname{iso}}
\newcommand{\cN}{\mathcal{N}}
\newcommand{\cO}{\mathcal{O}}
\newcommand{\cE}{\mathcal{E}}
\newcommand{\cK}{\mathcal{K}}
\newcommand{\cP}{\mathcal{P}}
\newcommand{\cQ}{\mathcal{Q}}

\newcommand{\footremember}[2]{%
  \footnote{#2}%
  \newcounter{#1}%
  \setcounter{#1}{\value{footnote}}%
}

\newcommand{\ag}[1]{\textcolor{purple}{(Anna: #1)}}
\newcommand{\sd}[1]{\textcolor{red}{(Sahar: #1)}}
\title{The number of isolated vertices at the two ends of the window: at least one and at most d below, none above}
\author{Sahar Diskin, Anna Geisler}
\date{arxiv:2404.14020v4; CC BY 4.0}
\begin{document}
\maketitle
\noindent\textit{Editorial note.} All delivered mathematical text is the author's own, copied from the pinned author source sources/combinatorics-fl04-vspace-1-5cm-perfect-matching-in-p/pinned-original-source.tex; no statement or proof step was rewritten, reordered or removed.\par\medskip
\noindent\textit{Editorial note.} Preamble provenance: the source pins the publisher class scrartcl, which the delivery wrapper does not accept, so this card uses sources/combinatorics-fl04-vspace-1-5cm-perfect-matching-in-p/retained-preamble.tex, which supplies the amsart baseline and then carries the author's own newtheorem and macro declarations (pinned source lines 30-57) verbatim. No macro was invented and no mathematical symbol was changed.\par\medskip
\noindent\textit{Editorial note.} Bibliography provenance: the source carries no thebibliography environment and ends with bibliographystyle abbrv plus bibliography perc (author source lines 1691-1692), so the bibliography is resolved with BibTeX from the identical copy staged with the same arXiv version under sources/product-graph-matchings-2404.14020v4/perc.bib; no entry was edited and the author's own abbrv style is used. The counted proof itself contains no citation.\par\medskip
\noindent\textit{Editorial note.} Selection provenance: this card was opened under the file-based control contract (CONTROL.md), with the mechanical pre-card gate run on its statement and proof interval before the card was written and reporting the interval free.\par\medskip
\noindent\textit{Editorial note.} External references are navigation pointers only; every deliberately kept step of the author's argument stays inside the counted spans.\par\medskip
\noindent\textit{Allowed standard prerequisites (editorial).}\begin{itemize}
\item The notation iso(G\_p) for the number of isolated vertices of the random subgraph, and the two ends p\_- and p\_+ of the window together with the parameter omega. Set up in the author's section on the random graph process (source lines 440-520); the expectation n(1-p)\textasciicircum{}d is computed at the start of the counted proof. Fixes the random variables whose size the statement bounds and the two probabilities at which they are evaluated.
\end{itemize}\medskip
\begin{lemma}\label{lem:number-isolates}
Whp,
\[
  1\le \iso(G_{p_-})\le d
  \qquad\text{and}\qquad
  \iso(G_{p_+})=0.
\]
\end{lemma}

\begin{proof}
Let $Y_p=\iso(G_p)$.
Then $\mathbb{E}Y_p=n(1-p)^d$.
For $p=p_+$ this expectation is $1/\omega$, and Markov's inequality gives $Y_{p_+}=0$ \whp.
For $p=p_-$, we have $\mathbb{E}Y_{p_-}=\omega$, and
\[
  \mathbb{P}(Y_{p_-}>d)\le\frac{\omega}{d}=o(1).
\]

For the lower bound, expand the second moment.
Non-adjacent pairs of vertices are simultaneously isolated with probability $(\omega/n)^2$. An adjacent pair of vertices has probability at most $\alpha_M(\omega/n)^2$, where $\alpha_M$ is a suitable constant.
Since $G$ has $nd/2$ edges,
\[
  \mathbb{E}Y_{p_-}^2
  \le \omega+\omega^2+O\left(\frac{d\omega^2}{n}\right)
  =\omega^2+O(\omega).
\]
Hence $\operatorname{Var}(Y_{p_-})=O(\omega)=o(\omega^2)$, and Chebyshev's inequality gives $Y_{p_-}>0$ \whp.
\end{proof}

\bibliographystyle{plain}
\bibliography{perc}
\end{document}
